\documentclass[10pt,a4paper]{amsart}
\usepackage[margin=19mm]{geometry}
\usepackage{amsmath,amssymb,mathtools}
\usepackage[scr=rsfs,cal=euler]{mathalfa}
\usepackage{xfrac}
\usepackage[unicode,hidelinks,bookmarks=false]{hyperref}
\newcommand{\ie}{i.e.\ }
\newcommand{\cf}{cf.\ }
\newcommand{\resp}{respectively\ }
\newcommand{\Subsection}{\subsection}
\providecommand{\nobreakdash}{\nobreak\hbox{-}}
\setlength{\emergencystretch}{2em}
\multlinegap0pt
\usepackage{bm}
\usepackage{bbm}
\usepackage{dsfont}
\usepackage{xspace}
\usepackage{cancel}
\usepackage{mathtools}
\usepackage{amsthm}
\makeatletter
\@ifundefined{theorem}{\newtheorem{theorem}{Theorem}}{}
\@ifundefined{lemma}{\newtheorem{lemma}[theorem]{Lemma}}{}
\makeatother
\title[Bottom Spectrum Estimate Under Curvature Integrability Condi]{Bottom Spectrum Estimate Under Curvature Integrability Condition}
\author{Cole Durham}
\date{Source: arXiv:2506.23997v1 (CC BY 4.0)}
\begin{document}
\maketitle

\noindent\emph{Source and licence.} Bottom Spectrum Estimate Under Curvature Integrability Condition. Version arXiv:2506.23997v1 (CC BY 4.0), distributed under the Creative Commons Attribution 4.0 International licence, as recorded in the arXiv licence metadata for this exact version and shown on the abstract page https://arxiv.org/abs/2506.23997v1. Full rights notice: licenses/arxiv-2506.23997-CC-BY-4.0.txt.\par\smallskip

\noindent\textit{Editorial scope.} Counted unit: the Lemma whose source label is \texttt{None} in the article (source0.tex lines 212--218, source label \texttt{None}), delivered together with the article's own complete original proof of it (source0.tex lines 220--330, one \texttt{proof} environment of 111 lines). No mathematics is written, repaired or re-derived: statement and proof are the article's own text line for line. Required context retained uncounted: none. Every definition, display and label of those lines is retained unchanged. Omitted from this package and not redistributed: 1--211, 219--219, 331--748. Nothing inside a retained excerpt is omitted or altered: every retained body is delivered exactly as the article's lines are written, so no omission receipt of any kind accompanies this package. Citations: the retained excerpts carry no citation command at all, so no bibliography entry of the article is transcribed and no reference is redistributed. Reference adaptation: none was needed and none was made; every reference inside the delivered unit resolves inside it, so no unresolved reference or citation is produced. Printed slips kept literal: no printed word, symbol, hypothesis or number of the counted statement or of its proof was altered. Layout and class: the article's own class is \texttt{amsart} with packages including ; the wrapper is the lane's pinned \texttt{amsart} editorial wrapper at 10pt on A4, which restates the theorem-like environments the retained text needs and adds the article's own macro definitions that the retained text needs (none), with extra wrapper packages (bm, bbm, dsfont, xspace, cancel, mathtools). The delivered numbering is the unit's own. Assets: no retained span contains a figure environment, a graphics-inclusion command, a PDF-page inclusion, a table or a TikZ picture; no image file is copied into this package, the package declares no asset dependency and no image file is redistributed with it. Licence: version arXiv:2506.23997v1 of this article is distributed under the Creative Commons Attribution 4.0 International licence, as recorded in the arXiv licence metadata for this exact version and shown on the abstract page https://arxiv.org/abs/2506.23997v1. A full rights notice is delivered at licenses/arxiv-2506.23997-CC-BY-4.0.txt. Characters: every retained excerpt is pure ASCII, so the delivered engine, XeLaTeX with its default Latin Modern text font, reports no missing character.
\begin{lemma}
    Let $M$ be a complete $n$-dimensional manifold with $\lambda_0(M)>0$ and suppose $w\in L^s(M)$ for some $s\geq \frac{3}{2}$. Then the minimal positive Green's function $G$ with pole at $p\in M$ satisfies
    \begin{equation*}
    \int_{M\setminus B_p(1)} \frac{|\nabla G|^3}{G^2\ln^{2\gamma}(1+G^{-1})}<\infty
    \end{equation*}
    for all $\gamma > \frac{1}{2}$.
\end{lemma}

\begin{proof}

   Let $v = \ln(G)$. Our goal is to estimate $\int_M |\nabla v|^3\phi^2$ where $\phi$ is a suitably chosen cutoff function which vanishes on $B_p(1)$. A standard manipulation of the Bochner formula combined with the definition of $w$ gives 
\begin{align*}
\frac{1}{n-1}|\nabla v|^3\leq& \Delta |\nabla v|-\frac{1}{n-1}|\nabla |\nabla v||^2|\nabla v|^{-1}+\frac{n-2}{n-1}\left<\nabla |\nabla v|^2,\nabla v\right>|\nabla v|^{-1}\\
&+w|\nabla v|+(n-1)|\nabla v|.
\end{align*}
From this inequality, for any cutoff function $\phi$ we have
\begin{align*}
\frac{1}{n-1}\int_M |\nabla v|^3\phi^2 \leq &\int_M (\Delta |\nabla v|)\phi^2 - \frac{1}{n-1}\int_M |\nabla |\nabla v||^2|\nabla v|^{-1}\phi^2\\
&+\frac{n-2}{n-1}\int_M\left<\nabla |\nabla v|^2,\nabla v\right>|\nabla v|^{-1}\phi^2+\int_M w|\nabla v|\phi^2\\
&+(n-1)\int_M |\nabla v|\phi^2. 
\end{align*}
Using integration by parts, this is equivalent to
\begin{align*}
\frac{1}{n-1}\int_M |\nabla v|^3\phi^2\leq &-\int_M \left<\nabla \phi^2,\nabla |\nabla v|\right>- \frac{1}{n-1}\int_M |\nabla |\nabla v||^2|\nabla v|^{-1}\phi^2 \\
&+\frac{2(n-2)}{n-1}\int_M\left<\nabla |\nabla v|,\nabla v\right>\phi^2 +\int_M w|\nabla v|\phi^2\\
&+(n-1)\int_M|\nabla v|\phi^2.
\end{align*}

Now we choose $\phi = G^{\frac{1}{2}}\eta$ for some cutoff function $0\leq \eta\leq 1$ such that $\eta = 0$ on $B_p(1)\cup (M\setminus B_p(2R))$. Expanding the right side of the previous inequality, we arrive at
\begin{align*}
\frac{1}{n-1}\int_M |\nabla v|^3\phi^2\leq &- \int_M \left<\nabla G, \nabla|\nabla v|\right>\eta^2-2\int_M \left<\nabla \eta, \nabla |\nabla v|\right>G\eta \\
&-\frac{1}{n-1}\int_M |\nabla |\nabla v||^2|\nabla v|^{-1}G\eta^2 + (n-1)\int_M|\nabla v|G\eta^2\\
&+\int_M w|\nabla v| G\eta^2  +\frac{2(n-2)}{n-1}\int_M \left<\nabla G, \nabla|\nabla v|\right> \eta^2.
\end{align*}
Integration by parts and the fact that $G$ is harmonic away from $p$ show
\begin{align}\label{eq1}
\frac{1}{n-1}\int_M |\nabla v|^3\phi^2 \leq &\left(1-\frac{2(n-2)}{n-1}\right)\int_M \left<\nabla G,\nabla\eta^2\right>|\nabla v|-2\int_M\left<\nabla\eta,\nabla |\nabla v|\right>G\eta \notag \\
&-\frac{1}{n-1}\int_M|\nabla |\nabla v||^2|\nabla v|^{-1}G\eta^2+\int_M w|\nabla v| G\eta^2\notag\\
&+ (n-1)\int_M|\nabla v|G\eta^2.  
\end{align}
We apply Cauchy's inequality with $\delta>0$ to each of the first two terms in (\ref{eq1}). First,
\begin{align*}
\left(1-\frac{2(n-2)}{n-1}\right)\int_M \left<\nabla G,\nabla\eta^2\right>|\nabla v|&\leq C\int_M \eta |\nabla \eta||\nabla v|^2 G\\
&\leq C\delta \int_M |\nabla v|^3G\eta^2 + \frac{C}{\delta}\int_M |\nabla v||\nabla \eta|^2 G.
\end{align*}
By similar reasoning
\begin{align*}
-2\int_M\left<\nabla\eta,\nabla |\nabla v|\right>G\eta &\leq C\int_M |\nabla\eta||\nabla|\nabla v|| G\eta \\
&\leq C\delta \int_M |\nabla|\nabla v||^2|\nabla v|^{-1}G\eta^2 +\frac{C}{\delta}\int_M|\nabla\eta|^2G|\nabla v|.
\end{align*}
We replace terms in (\ref{eq1}) by their respective bounds. Under the assumption $\delta\leq \frac{C}{n-1}$, we simplify and have
\begin{align}
\frac{1}{n-1}\int_M|\nabla v|^3\phi^2 \leq& C\delta\int_M |\nabla v|^3G\eta^2 + \left(C\delta - \frac{1}{n-1}\right)\int_M|\nabla |\nabla v||^2|\nabla v|^{-1}G\eta^2 \notag \\
&+ \frac{C}{\delta}\int_M (\eta^2+|\nabla\eta|^2)G|\nabla v|+\int_Mw|\nabla v|G\eta^2.\notag
\end{align}
Assuming $\delta$ is small enough such that $\frac{1}{n-1} - C\delta > 0$, given $|\nabla|\nabla v||^2|\nabla v|^{-1} G\eta^2 \geq 0$ we find
\begin{equation}\label{eq2}\left(\frac{1}{n-1}-C\delta\right)\int_M |\nabla v|^3\phi^2 \leq \frac{C}{\delta}\int_M(\eta^2+|\nabla\eta|^2)G|\nabla v|+\int_M w|\nabla v|G\eta^2.\end{equation}

To continue toward achieving a bound, we specify $\eta = \tau f(G)$ where $0\leq \tau \leq 1$ is a cutoff function satisfying $\tau = 0$ on $B_p(1)\cup (M\setminus B_p(2R)),$ $\tau = 1$ on $B_p(R)\setminus B_p(2)$. As in Lemma 2.1,
\begin{equation*}
f(G) = \frac{1}{\ln^\gamma(AG^{-1})}, \;\; A = e^{\frac{1}{\alpha}} \max_{\partial B_p(1)} G,\;\; \gamma>1/2,
\end{equation*}
but here we require $0<\alpha<\delta$ is small relative to $\delta$. By Young's inequality for $0<\beta<\delta$ we know
\begin{align}\label{eq3}
\int_M |\nabla v|G\eta ^2&\leq \frac{\beta}{3}\int_M\frac{|\nabla G|^3}{G^2}f^2\tau^2+\frac{2}{3\beta^\frac{1}{2}}\int_M Gf^2\tau^2 \notag\\
&\leq \frac{\beta}{3}\int_M |\nabla v|^3\phi^2+\frac{C}{\beta^\frac{1}{2}}
\end{align}
where we have used Lemma 2.1 in the second inequality. Next, expanding the gradient of $\eta$ shows
\begin{equation*}
\int_M |\nabla \eta|^2G|\nabla v|\leq 2\int_M|\nabla \tau|^2f^2|\nabla G|+2\int_M\tau^2|\nabla f|^2|\nabla G|.
\end{equation*}
Since $\tau$ is assumed to be constant everywhere except on the concentric annuli $(B_p(2)\setminus B_p(1))\cup (B_p(2R)\setminus B_p(R))$, we use the Cauchy-Schwarz inequality, Lemma 2.1, and the co-area formula to write
\begin{align*}
\int_M |\nabla \tau|^2f^2|\nabla G|\ 
&\leq \int_{M\setminus B_p(1)}\left(\frac{|\nabla G|}{G^\frac{1}{2}}f\right)\left(G^\frac{1}{2}f\right)\\
& \leq \left(\int_{L(0,e^{-\frac{1}{\alpha}}A)}\frac{|\nabla G|^2}{G\ln^{2\gamma}(AG^{-1})}\right)^\frac{1}{2}\left(\int_{M\setminus B_p(1)}Gf^2\right)^\frac{1}{2}\\
& \leq C\left(\int_0^{e^{-\frac{1}{\alpha}}A}\frac{dr}{r\ln^{2\gamma}(A/r)}\right)^\frac{1}{2}\\
& \leq C.
\end{align*}
From calculations in Lemma 2.1, recall
\begin{equation*}
|\nabla f| = \gamma \frac{|\nabla G|}{G\ln^{\gamma+1}(AG^{-1})},
\end{equation*}
from which it follows 
\begin{equation*}
\tau^2|\nabla f|^2|\nabla G| = \frac{\gamma^2}{\ln^2(AG^{-1})}|\nabla v|^3\phi^2.
\end{equation*}
From the definition of $A = e^\frac{1}{\alpha}\max_{\partial B_p(1)}G$ we know that $\ln^2(AG^{-1})\geq \frac{1}{\alpha^2}$ on $M\setminus B_p(1)$, and hence
\begin{equation*}
\int_M |\nabla \eta|^2G|\nabla v|\leq 2\int_M |\nabla\tau|^2f^2|\nabla G| +2\gamma^2\alpha^2\int_M |\nabla v|^3\phi^2 .
\end{equation*}
In (\ref{eq2}), these estimates show
\begin{equation*}
\left(\frac{1}{n-1}-C\delta-\frac{C}{\delta}\left(\gamma^2\alpha^2+\beta\right)\right)\int_M|\nabla v|^3\phi^2\leq \frac{C}{\delta}\left(1+\frac{1}{\beta^\frac{1}{2}}\right)+\int_M w|\nabla v| \phi^2.
\end{equation*}

We lastly turn our attention to the curvature term. Since our hypotheses state $w\in L^s(M)$ for some $s\geq\frac{3}{2}$, an index $t$ conjugate to $s$ satisfies $1\leq t\leq 3$ and therefore
\begin{align*}
\int_M w|\nabla v|\phi^2 &\leq \frac{1}{s\delta^\frac{s}{t}}\int_M w^s\phi^2+\frac{\delta}{t}\int_M |\nabla v|^t\phi^2\\
&\leq \frac{C}{\delta^\frac{s}{t}}+\frac{\delta}{t}\int_M|\nabla v|^t\phi^2.
\end{align*}
Once more applying Young's inequality, we see
\begin{equation*}
\int_M |\nabla v|^t\phi^2\leq \left(\frac{t-1}{2}\right)\int_M |\nabla v|^3+\left(\frac{3-t}{2}\right)\int_M|\nabla v|\phi^2.
\end{equation*}
Then by (\ref{eq3}), since $\frac{t-1}{2}, \frac{3-t}{2}\leq 1$ for $1\leq t\leq 3$,
\begin{align*}
\frac{\delta}{t}\int_M |\nabla v|^t\phi^2 &\leq \delta\int_M |\nabla v|^3+\delta\left\{\beta\int_M |\nabla v|^3\phi^2+\frac{C}{\beta^\frac{1}{2}}\right\}
\end{align*}
which implies
\begin{equation*}
\int_M w|\nabla v|\phi^2\leq \frac{C}{\delta^\frac{s}{t}} + C\delta \int_M|\nabla v|^3\phi^2 + C\frac{\delta}{\beta^\frac{1}{2}}.
\end{equation*}
We may then conclude
\begin{equation*}
\left(\frac{1}{n-1}-C\delta -\frac{C}{\delta}(\gamma^2\alpha^2+\beta)\right)\int_M|\nabla v|^3\phi^2\leq \frac{C}{\delta}\left(1+\frac{1}{\beta^\frac{1}{2}}\right)+\frac{C}{\delta^\frac{s}{t}}.
\end{equation*}
The parameters $\alpha,\beta,$ and $\delta$ can be chosen such that $\frac{1}{n-1}-C\delta -\frac{C}{\delta}(\gamma^2\alpha^2+\beta)>0$, thus letting $R\to\infty$ in the definition of $\phi$ proves the result.
\end{proof}

\end{document}
